\section{Measurable Realization and Dichotomy}
\label{sec:measurable-realization-and-dichotomy}

The sample space is \(\R^{\N}\), not \(\ell^2\), so a bounded operator on
\(\ell^2\) does not act pointwise by a naively convergent matrix product.
This section constructs the required measurable versions.
Then it identifies the global measure class through likelihood-ratio martingales
and a dense translation group.

\subsection{Measurable linear actions and countable-orbit realizations}
\label{sec:measurable-realization}

Throughout this subsection, only \(\E X=0\) and \(\E X^2=1\) for \(X\sim\rho\) are needed.
Recall \(\mu=\rho^{\otimes\N}\) and \(\mu_n=\rho^{\otimes n}\).
The sample sequence is not an \(\ell^2\)-valued random variable.
So we first specify how an operator on \(\ell^2\) acts on its law.

\begin{lemma}[Measurable realization on each countable operator orbit]
\label{lem:measurable-action}
For every \(T\in\mathcal B(H)\), there is a Borel map
\(\Phi_T:\R^{\N}\to\R^{\N}\) whose coordinates are the \(L^2(\mu)\) row sums
\[
  (\Phi_Tx)_i=\sum_{j\geq1}T_{ij}x_j.
\]
These maps can be chosen so that the following statements hold.
If a probability measure \(\lambda\) satisfies
\begin{equation}
  \int\left|\sum_j a_jx_j\right|^2\,d\lambda(x)
  \leq C_\lambda\|a\|_2^2
  \quad\text{for every finitely supported }a,
  \label{eq:covariance-domination}
\end{equation}
then \((\Phi_T)_\#\lambda\) also satisfies this condition.
Its constant is at most \(C_\lambda\|T\|_{\mathrm{op}}^2\).  Also
\begin{equation}
  \Phi_S\circ\Phi_T=\Phi_{ST}
  \quad\lambda\text{-almost surely}
  \qquad(S,T\in\mathcal B(H)).
  \label{eq:measurable-composition}
\end{equation}
For each fixed countable subgroup \(\Gamma\leq\GL(H)\),
there is a \(\Gamma\)-invariant Borel set \(D\subset\R^{\N}\) of full
measure for every \(\lambda\) satisfying
\eqref{eq:covariance-domination}.
On this set, these restrictions form an action by bimeasurable
automorphisms.
In particular, this applies to \(\mu\) and all its \(\Gamma\)-orbit laws,
including singular orbit laws.
\end{lemma}

\begin{proof}
For \(a\in\ell^2\), the finite sums \(\sum_{j\leq n}a_jx_j\) have a
unique \(L^2(\lambda)\) limit \(L_a\), by
\eqref{eq:covariance-domination}.  Choose
\[
  n_k(a)=\min\left\{n\geq k:
                \sum_{j>n}|a_j|^2\leq2^{-4k}\right\}.
\]
Define \(\widehat L_a(x)\) to be the limit of
\(\sum_{j\leq n_k(a)}a_jx_j\) when it exists, and zero otherwise.
The estimate
\[
  \left\|\sum_{j\leq n_k(a)}a_jx_j-L_a\right\|_{L^2(\lambda)}^2
  \leq C_\lambda2^{-4k}
\]
and Borel--Cantelli show that \(\widehat L_a=L_a\),
\(\lambda\)-almost surely.  This one definition therefore works for every
law satisfying \eqref{eq:covariance-domination}; it does not depend on
knowing whether that law is equivalent to \(\mu\).  Set
\(\Phi_T(x)_i=\widehat L_{T^*e_i}(x)\).  The construction is Borel,
and the same is true jointly in the matrix coefficients of \(T\) and in
\(x\), since each tail norm and the integer choice \(n_k\) is Borel.

For finitely supported \(b\), the finite sum of output coordinates equals
\(L_{T^*b}\) almost surely.  This proves the asserted second-moment bound
for the pushforward.  For the \(i\)-th row of \(S\), finite output sums
satisfy
\[
  \sum_{j\leq n}S_{ij}\widehat L_{T^*e_j}
     =L_{T^*P_nS^*e_i}
     \longrightarrow L_{T^*S^*e_i}
     \quad\text{in }L^2(\lambda).
\]
The selected subsequence on the left also converges almost surely to
\(\Phi_S(\Phi_T(x))_i\), because the output law has the required
second-moment bound.  Uniqueness of the limit proves
\eqref{eq:measurable-composition}.

For a countable \(\Gamma\), let \(B\) be the Borel set on which all its
row limits, identity relations, and pairwise composition relations hold.
The preceding argument gives \(\lambda(B)=1\) for every law under
consideration.  Put
\[
  D=\bigcap_{T\in\Gamma}\Phi_T^{-1}(B).
\]
Every orbit law satisfies the same second-moment condition, so
\(\lambda(D)=1\).  The composition identities imply that
\(D\) is invariant and that \(\Phi_{T^{-1}}\) is the inverse of
\(\Phi_T\) there.  This proves the final assertion.
\end{proof}

Write \(\nu_T=(\Phi_T)_\#\mu\).
The preceding lemma makes assertions about a common carrier only for a
fixed countable subgroup.
No single pointwise action of all of \(\GL(H)\) on a common
conull set is required.
The subsequences in its proof are also essential.
A covariance bound gives convergence of arbitrary row sums in \(L^2\).
That is not enough for almost-sure convergence of all ordinary partial sums
under a correlated orbit law.
On the original independent product law, the versions agree with the usual
almost-sure row sums.

\subsection{Rational translation invariance and the zero--one dichotomy}
\label{sec:translation-dichotomy}

\begin{lemma}[Dichotomy for every bounded invertible orbit]
\label{lem:orbit-dichotomy}
For every \(T\in\GL(H)\),
\begin{equation}
  \nu_T\sim\mu\quad\text{or}\quad\nu_T\perp\mu.
  \label{eq:orbit-dichotomy}
\end{equation}
This conclusion uses only positivity of \(f\), its centered finite second
moment, and \(\sqrt f\in W^{1,2}(\R)\).
In particular, it holds under \eqref{eq:main-assumptions}.
\end{lemma}

\begin{proof}
Put \(g=\sqrt f\).  The translation estimate
\[
  1-\Aff(f,f(\cdot-a))
     =\tfrac12\|g(\cdot-a)-g\|_2^2
     \leq\tfrac12\|g'\|_2^2a^2
\]
and Kakutani's theorem imply
\((x\mapsto x+a)_\#\mu\sim\mu\) for every \(a\in\ell^2\).
The row realizations satisfy
\[
  \Phi_T(x+T^{-1}a)=\Phi_T(x)+a
  \quad\mu\text{-almost surely};
\]
this follows directly by adding the absolutely convergent scalar row
pairing with the \(\ell^2\) vector \(T^{-1}a\).
Consequently, \(\nu_T\) is also quasi-invariant under every
\(\ell^2\) translation.

Let \(\mathcal A\) be the countable additive group generated by
\(\{qe_j,qTe_j:q\in\mathbb Q,\ j\geq1\}\).
Both \(\mu\) and \(\nu_T\) are quasi-invariant and ergodic under
translation by \(\mathcal A\).  Indeed, an event invariant modulo \(\mu\) under
all rational finite-coordinate translations has, by positivity of the
finite-dimensional densities, almost-everywhere constant sections in each
finite coordinate block.  It is therefore a tail event modulo \(\mu\),
and has probability zero or one.  Here one uses the elementary fact that
a measurable function on \(\R^n\) invariant almost everywhere
under rational translations is constant almost everywhere, obtained, for
example, by convolution with smooth approximate identities.
For \(\nu_T\), pull an invariant event back by \(\Phi_T\); invariance
under \(qTe_j\) becomes invariance under \(qe_j\), giving the same
zero--one conclusion.

Write \(\nu_T=\nu_a+\nu_s\) for its Lebesgue decomposition relative to
\(\mu\).  Choose a Borel \(\mu\)-null set \(N_0\) supporting
\(\nu_s\), and let \(N=\bigcup_{a\in\mathcal A}(N_0+a)\).
Quasi-invariance of \(\mu\) makes \(N\) a \(\mu\)-null set;
it is exactly \(\mathcal A\)-invariant and still supports \(\nu_s\).
Thus \(\nu_T(N)=\nu_s(\R^{\N})\), and ergodicity of \(\nu_T\) makes this
number zero or one.  The latter case gives singularity.  In the former
case, write \(\nu_T=Z\mu\).  Quasi-invariance of both measures implies
that \(\{Z>0\}\) is invariant modulo \(\mu\) under \(\mathcal A\).
Its \(\mu\)-measure is positive because \(\int Z\,d\mu=1\), and
ergodicity of \(\mu\) makes that measure one.  This gives equivalence.
\end{proof}

\subsection{Finite-marginal likelihood ratios and Hellinger energy}
\label{sec:finite-marginal-likelihood-ratios}

\begin{lemma}[Finite-marginal likelihood ratios and two-sided uniform integrability]
\label{lem:finite-marginal-likelihood-ratios}
Assume \(f>0\) almost everywhere, \(\sqrt f\in W^{1,2}(\R)\),
and \(\E X=0\), \(\E X^2=1\).
For \(T\in\GL(H)\), every finite output law
\((\pi_n)_\#\nu_T\) is equivalent to \(\mu_n\).
Writing \(\mathcal F_n=\sigma(x_1,\ldots,x_n)\), define
\begin{equation}
 \begin{aligned}
  \ell_n(x)=\frac{d[(\pi_n)_\#\nu_T]}{d\mu_n}(\pi_n x),
  &\qquad E_n(T)=-2\log\E_\mu\sqrt{\ell_n},\\
  \Aff\bigl(\mu_n,(\pi_n)_\#\nu_T\bigr)
  &=\E_\mu\sqrt{\ell_n}.
 \end{aligned}
  \label{eq:finite-marginal-likelihood-ratios}
\end{equation}
Then \((\ell_n)\) is a nonnegative mean-one \(\mu\)-martingale.
We have
\begin{equation}
  \Aff\bigl(\mu_n,(\pi_n)_\#\nu_T\bigr)\downarrow\Aff(\mu,\nu_T),
  \qquad
  \nu_T\sim\mu
    \Longleftrightarrow\sup_nE_n(T)<\infty.
  \label{eq:marginal-affinity-limit}
\end{equation}
Otherwise \(E_n(T)\uparrow\infty\).
Then \(\nu_T\perp\mu\).
If \(\nu_T\sim\mu\) with density \(Z_T\), then
\begin{align}
  \ell_n&=\E_\mu[Z_T\mid\mathcal F_n]
       \longrightarrow Z_T
       &&\text{in }L^1(\mu),
       \label{eq:forward-likelihood-ui}\\
  \ell_n^{-1}&=\E_{\nu_T}[Z_T^{-1}\mid\mathcal F_n]
       \longrightarrow Z_T^{-1}
       &&\text{in }L^1(\nu_T).
       \label{eq:reverse-likelihood-ui}
\end{align}
In particular, both likelihood-ratio martingales are uniformly integrable in
their respective reference measures.
\end{lemma}

\begin{proof}
The first \(n\) rows of \(T\) are linearly independent, since \(T^*\)
is injective.  Some \(n\) input columns therefore form an invertible
\(n\times n\) minor \(A\).  Split the input into these coordinates
and their complement.  The first \(n\) outputs have the form
\(AU+W\), where \(U\sim\mu_n\) and \(W\) is independent of \(U\).
Conditioning on \(W\) writes their density as a mixture of translates
of the strictly positive density of \(AU\).  This mixture is positive
and finite almost everywhere, proving the finite-dimensional equivalence.

Consistency of projected laws proves the martingale assertion.  The
martingale limit is
\(\ell_\infty=d(\nu_T)_a/d\mu\), the density of the absolutely continuous
part of \(\nu_T\).  To see the identification explicitly, Fatou's lemma
on cylinder sets gives \(\ell_\infty\mu\leq\nu_T\), first on the cylinder
algebra and then on the generated sigma-field.  Conversely, if
\((\nu_T)_a=Z\mu\), then
\(\E_\mu[Z\mid\mathcal F_n]\leq \ell_n\), so martingale
convergence gives \(Z\leq \ell_\infty\).  The two inequalities identify
the limit.
Since \(\E_\mu(\sqrt{\ell_n})^2=1\), the square roots are uniformly
integrable in \(L^1(\mu)\).  They thus converge in \(L^1\) to
\(\sqrt{\ell_\infty}\).  Conditional Jensen's inequality gives the
monotonicity of their expectations, and consequently
\[
  \lim_n\Aff\bigl(\mu_n,(\pi_n)_\#\nu_T\bigr)
      =\int\sqrt{\ell_\infty}\,d\mu
      =\Aff(\mu,\nu_T).
\]
The dichotomy in \cref{lem:orbit-dichotomy} now proves
\eqref{eq:marginal-affinity-limit} and the singular alternative.

Under equivalence, the forward conditional-expectation identity is the
definition of a finite-marginal density.  The reverse identity follows because
for \(B\in\mathcal F_n\),
\[
  \int_B \ell_n^{-1}\,d\nu_T=\mu(B)
       =\int_B Z_T^{-1}\,d\nu_T.
\]
Conditional-expectation convergence gives both \(L^1\) limits and hence
both uniform-integrability assertions.
\end{proof}

By data processing, the supremum of the projected Hellinger energies is
unchanged if one takes all finite coordinate marginals: the initial
segments are cofinal among finite coordinate sets.  Thus the criterion
does not depend on the enumeration of the coordinates.
